%%%%%%%%%%%%%%%%%%%%%%%%%%%%%%%%%%%%%%%%%%%%%%%%%%%%%%%%%%%%%%%%
%
%  Template for homework of Machine Learning.
%%%%%%%%%%%%%%%%%%%%%%%%%%%%%%%%%%%%%%%%%%%%%%%%%%%%%%%%%%%%%%%%

\documentclass[11pt,letter,notitlepage]{ctexart}

\usepackage{fontspec}
\setCJKmainfont{KaiTi}[BoldFont=KaiTi,ItalicFont=KaiTi]

% Mise en page
\usepackage[left=2cm, right=2cm, lines=45, top=0.8in, bottom=0.7in]{geometry}
\usepackage{fancyhdr}
\usepackage{fancybox}
\usepackage{graphicx}
\usepackage{pdfpages}
\usepackage{enumitem}
\usepackage{algorithm}
\usepackage{algorithmic}
\usepackage{array}
\usepackage{booktabs}
\renewcommand{\headrulewidth}{1.5pt}
\renewcommand{\footrulewidth}{1.5pt}
\newcommand\Loadedframemethod{TikZ}
\usepackage[framemethod=\Loadedframemethod]{mdframed}
\usepackage{amssymb,amsmath}
\usepackage{amsthm}
\usepackage{thmtools}
\usepackage{bm}
\usepackage{multirow}
\usepackage{makecell}
\usepackage{tabularx}
\usepackage{longtable}
\usepackage{float}

\setlength{\topmargin}{0pt}
\setlength{\textheight}{9in}
\setlength{\headheight}{0pt}
\setlength{\oddsidemargin}{0.25in}
\setlength{\textwidth}{6in}

%%%%%%%%%%%%%%%%%%%%%%%%
%%%%%% Define math operator %%%%%
%%%%%%%%%%%%%%%%%%%%%%%%
\DeclareMathOperator*{\argmin}{\bf argmin}
\DeclareMathOperator*{\relint}{\bf relint\,}
\DeclareMathOperator*{\dom}{\bf dom\,}
\DeclareMathOperator*{\intp}{\bf int\,}

\DeclareMathOperator*{\cl}{\bf cl\,}
\DeclareMathOperator*{\bd}{\bf bd\,}
\DeclareMathOperator*{\conv}{\bf conv\,}
\DeclareMathOperator*{\epi}{\bf epi\,}

%%%%%%%%%%%%%%%%%%%%%%%
\pagestyle{fancy}

%%%%%%%%%%%%%%%%%%%%%%%%
%% Define the Exercise environment %%
%%%%%%%%%%%%%%%%%%%%%%%%
\mdtheorem[
topline=false,
rightline=false,
leftline=false,
bottomline=false,
leftmargin=-10,
rightmargin=-10
]{exercise}{\textbf{Exercise}}

%%%%%%%%%%%%%%%%%%%%%%%%
%% Define the Problem environment %%
%%%%%%%%%%%%%%%%%%%%%%%%
\mdtheorem[
topline=false,
rightline=false,
leftline=false,
bottomline=false,
leftmargin=-10,
rightmargin=-10
]{problem}{\textbf{Problem}}

%%%%%%%%%%%%%%%%%%%%%%%
%% Define the Solution Environment %%
%%%%%%%%%%%%%%%%%%%%%%%
\declaretheoremstyle
[
spaceabove=0pt,
spacebelow=0pt,
headfont=\normalfont\bfseries,
notefont=\mdseries,
notebraces={(}{)},
headpunct={:},
headindent={},
postheadspace={ },
bodyfont=\normalfont,
qed=$\blacksquare$,
preheadhook={\begin{mdframed}[style=myframedstyle]},
postfoothook=\end{mdframed},
]{mystyle}

\declaretheorem[style=mystyle,title=Solution]{solution}

\mdfdefinestyle{myframedstyle}{%
    topline=false,
    rightline=false,
    leftline=false,
    bottomline=false,
    skipabove=-6ex,
    leftmargin=-10,
    rightmargin=-10,
    backgroundcolor=black!5}

%%%%%%%%%%%%%%%%%%%%%%%
%% Useful commands
%%%%%%%%%%%%%%%%%%%%%%%
\renewcommand{\eqref}[1]{Eq.~(\ref{#1})}
\newcommand{\figref}[1]{Fig.~\ref{#1}}
\newcommand{\proj}[2]{\textbf{P}_{#2} (#1)}
\newcommand{\lspan}[1]{\textbf{span} (#1)}
\newcommand{\rank}[1]{\textbf{rank} (#1)}
\newcommand{\tr}{\operatorname{tr}}
\newcommand{\RNum}[1]{\uppercase\expandafter{\romannumeral #1\relax}}
\newcommand{\mb}[1]{\mathbf{#1}}

% Definition environment
\theoremstyle{definition}
\newtheorem{definition}{Definition}

%%%%%%%%%%%%%%%%%%%%%%%%%%%%%%%%%%%%%%%%%%
%% Homework info
%%%%%%%%%%%%%%%%%%%%%%%%%%%%%%%%%%%%%%%%%%
\newcommand{\posted}{\text{2026.3.13}}
\newcommand{\due}{\text{2026.3.20}}
\newcommand{\hwno}{\text{1}}

%%%%%%%%%%%%%%%%%%%%%%%%%%%%%%%%%%%%
%% Put your information here
%%%%%%%%%%%%%%%%%%%%%%%%%%%%%%%%%%%%
\newcommand{\lec}{\text{肖力}}
\newcommand{\name}{\text{陈璟皓}}
\newcommand{\id}{\text{PB23010359}}

\lhead{\textbf{\name}}
\rhead{\textbf{\id}}
\chead{\textbf{Homework \hwno}}

\begin{document}
    \vspace*{-4\baselineskip}
    \thispagestyle{empty}

    \begin{center}
        {\bf\large 机器学习B}
    \end{center}

    \noindent
    Lecturer: \lec
    \hfill
    Homework \hwno
    \\
    Posted: \posted
    \hfill
    Due: \due
    \\
    Name: \name
    \hfill
    ID: \id

    \noindent
    \rule{\textwidth}{2pt}

    \medskip

    %%%%%%%%%%%%%%%%%%%%%%%%%%%%%%%%%%%%%%%%%%%%%%%%%%%%%%%%%%%%
    % Exercise 1
    %%%%%%%%%%%%%%%%%%%%%%%%%%%%%%%%%%%%%%%%%%%%%%%%%%%%%%%%%%%%
    \begin{exercise}[ROC曲线与AUC计算]
    假设给出如下表中 $10$ 个样本的真正类别和基于某个模型的它们预测为正类的概率, 绘制出 ROC 曲线, 并计算 AUC 值（需给出详细的计算过程）. 

    \begin{center}
        \renewcommand{\arraystretch}{1}
        \setlength{\tabcolsep}{30pt}
        \begin{tabular}{|c|c|c|}
            \hline
            样本 ID & 原本类别 & 预测为正类的概率 \\
            \hline
            1  & 正类 & 0.42 \\
            2  & 正类 & 0.65 \\
            3  & 负类 & 0.52 \\
            4  & 正类 & 0.86 \\
            5  & 正类 & 0.95 \\
            6  & 负类 & 0.53 \\
            7  & 负类 & 0.70 \\
            8  & 负类 & 0.43 \\
            9  & 正类 & 0.55 \\
            10 & 负类 & 0.35 \\
            \hline
        \end{tabular}
    \end{center}
    \end{exercise}

    \newpage

\begin{solution}

\end{solution}

    \newpage

    %%%%%%%%%%%%%%%%%%%%%%%%%%%%%%%%%%%%%%%%%%%%%%%%%%%%%%%%%%%%
    % Exercise 2
    %%%%%%%%%%%%%%%%%%%%%%%%%%%%%%%%%%%%%%%%%%%%%%%%%%%%%%%%%%%%
    \begin{exercise}[模型评估指标比较]
    在某二分类任务中, 现有两个模型 A 和 B, 在同一个测试集（共 $200$ 个样本）上的预测结果如下. 

    \medskip
    \textbf{模型 A 的混淆矩阵: }

    \begin{center}
    \renewcommand{\arraystretch}{1.3}
    \begin{tabular}{|c|c|c|}
        \hline
         & 预测正类 & 预测负类 \\
        \hline
        真实正类 & 48 & 12 \\
        \hline
        真实负类 & 18 & 122 \\
        \hline
    \end{tabular}
    \end{center}

    \medskip
    \textbf{模型 B 的混淆矩阵: }

    \begin{center}
    \renewcommand{\arraystretch}{1.3}
    \begin{tabular}{|c|c|c|}
        \hline
         & 预测正类 & 预测负类 \\
        \hline
        真实正类 & 54 & 6 \\
        \hline
        真实负类 & 30 & 110 \\
        \hline
    \end{tabular}
    \end{center}

    \medskip
    请完成以下问题: 
    \begin{enumerate}[label=(\arabic*)]
        \item 在 Accuracy、Precision、Recall、F1、FPR、FNR 指标上比较两个模型; 
        \item 若该任务是一般商品推荐系统（正类表示“用户感兴趣商品”, 负类表示“用户不感兴趣商品”）, 更适合选择哪个模型？为什么？
        \item 若该任务是重大疾病筛查系统（正类表示“患病”, 负类表示“未患病”）, 更适合选择哪个模型？为什么？
    \end{enumerate}
    \end{exercise}
    
    \newpage

\begin{solution}

\end{solution}

\end{document}